\documentclass[11pt,letterpaper]{article}

% ── packages ──────────────────────────────────────────────────────────────────
\usepackage[margin=1in]{geometry}
\usepackage{amsmath,amssymb}
\usepackage{graphicx}
\usepackage{booktabs}
\usepackage{multirow}
\usepackage{caption}
\usepackage{subcaption}
\usepackage{hyperref}
\usepackage{xcolor}
\usepackage{enumitem}
\usepackage{float}
\usepackage{microtype}
\usepackage[numbers,sort&compress]{natbib}

\hypersetup{colorlinks=true, linkcolor=blue!70!black, citecolor=blue!70!black, urlcolor=blue!70!black}

\captionsetup{font=small, labelfont=bf}
\setlength{\parskip}{4pt}

% ── macros ────────────────────────────────────────────────────────────────────
\newcommand{\method}[1]{\textbf{#1}}
\newcommand{\metric}[1]{\textit{#1}}

% ── title ─────────────────────────────────────────────────────────────────────
\title{%
  \Large\textbf{Cancer Subtype Discovery and Classification}\\[4pt]
  \large from TCGA Pan-Cancer Gene Expression Data\\[6pt]
  \normalsize 02-620 Machine Learning for Scientists --- Spring 2026
}
\author{Tolga Ozgun \and Yuchen Shao \and Roy Hung}
\date{April 2026}

\begin{document}
\maketitle

% ─────────────────────────────────────────────────────────────────────────────
\begin{abstract}
We present an end-to-end machine learning pipeline for cancer subtype discovery
and classification applied to bulk RNA-seq data from The Cancer Genome Atlas
(TCGA) Pan-Cancer dataset.  Starting from raw HTSeq-FPKM-UQ expression
matrices across six cancer types (BRCA, LUAD, KIRC, PRAD, COAD, GBM), we
implement three methods entirely from scratch: (1)~unsupervised clustering via
PCA + K-Means, (2)~supervised one-vs-rest logistic regression with
$\ell_2$-regularisation, and (3)~a multi-layer perceptron (MLP) implemented
in PyTorch.  All supervised models are evaluated using 5-fold stratified
cross-validation with preprocessing fit inside each fold to prevent data
leakage.  K-Means achieves an Adjusted Rand Index (ARI) of 0.794 and
Normalised Mutual Information (NMI) of 0.849, confirming that the six cancer
types are intrinsically well-separated in gene expression space.  Both
classifiers reach macro-F1 $\approx 0.993$ on the held-out test set, with
virtually no gap between the linear (LR) and non-linear (MLP) models --- a
finding consistent with the linear separability observed in the PCA projection.
\end{abstract}

% ─────────────────────────────────────────────────────────────────────────────
\section{Introduction}
\label{sec:intro}

Cancer is not a single disease.  Even within the same organ, tumours can differ
dramatically in molecular profile, prognosis, and response to therapy.  Bulk
RNA-sequencing captures a genome-wide snapshot of gene expression in a tumour
sample, providing a high-dimensional view of its molecular state.  The challenge
we address is threefold:

\begin{enumerate}[leftmargin=2em,itemsep=1pt]
  \item \textbf{Unsupervised discovery.}  Can meaningful cancer subtypes be
        recovered from unlabelled gene expression data using PCA followed by
        K-Means clustering?
  \item \textbf{Linear classification.}  How accurately can a multi-class
        logistic regression model (trained on PCA-reduced features) predict
        cancer type?
  \item \textbf{Non-linear classification.}  Does a deep neural network offer
        a measurable gain over the linear baseline?
\end{enumerate}

Identifying tumour subtypes from expression data has direct clinical
implications: patients with distinct molecular profiles often respond
differently to therapies~\cite{weinstein2013}.  The TCGA Pan-Cancer dataset
provides one of the largest and most carefully curated multi-cancer genomic
resources available~\cite{hutter2018}, making it an ideal testbed for
benchmarking ML methods on a high-stakes biological problem.

% ─────────────────────────────────────────────────────────────────────────────
\section{Dataset}
\label{sec:dataset}

\subsection{Source}

We use the TCGA Pan-Cancer (PANCAN) gene expression dataset downloaded via the
GDC Data Portal API\footnote{\url{https://portal.gdc.cancer.gov/}}.
Specifically, we download:

\begin{itemize}[leftmargin=2em,itemsep=1pt]
  \item \textbf{Expression matrix} (UUID
        \texttt{3586c0da-64d0-4b74-a449-5ff4d9136611}):
        \texttt{EBPlusPlusAdjustPANCAN\_IlluminaHiSeq\_RNASeqV2.geneExp.tsv}
        --- HTSeq-FPKM-UQ values, batch-corrected across all TCGA projects.
  \item \textbf{Clinical metadata} (UUID
        \texttt{1b5f413e-a8d1-4d10-92eb-7c4ae739ed81}):
        \texttt{TCGA-CDR-SupplementalTableS1.xlsx} --- patient-level cancer type
        annotations.
\end{itemize}

\subsection{Cancer Types Selected}

We focus on six cancer types with well-represented sample counts
(Table~\ref{tab:dataset}), chosen to span distinct tissue origins:
breast (BRCA), lung adenocarcinoma (LUAD), kidney renal clear cell (KIRC),
prostate (PRAD), colon (COAD), and glioblastoma multiforme (GBM).

\begin{table}[H]
  \centering
  \caption{Dataset composition after preprocessing.}
  \label{tab:dataset}
  \begin{tabular}{lrr}
    \toprule
    Cancer Type & Abbreviation & Samples \\
    \midrule
    Breast Invasive Carcinoma              & BRCA & 1{,}215 \\
    Kidney Renal Clear Cell Carcinoma      & KIRC &   606 \\
    Lung Adenocarcinoma                    & LUAD &   576 \\
    Prostate Adenocarcinoma                & PRAD &   550 \\
    Colon Adenocarcinoma                   & COAD &   492 \\
    Glioblastoma Multiforme                & GBM  &   168 \\
    \midrule
    \textbf{Total}                         &      & \textbf{3,607} \\
    \bottomrule
  \end{tabular}
\end{table}

Figure~\ref{fig:class_dist} shows the class distribution.  BRCA is the largest
class (1,215 samples) and GBM is the smallest (168 samples), introducing mild
class imbalance that we address by using macro-averaged metrics.

\begin{figure}[H]
  \centering
  \includegraphics[width=0.72\textwidth]{../results/report_figures/fig1_class_distribution.pdf}
  \caption{Sample count per cancer type in the full dataset.}
  \label{fig:class_dist}
\end{figure}

% ─────────────────────────────────────────────────────────────────────────────
\section{Data Preprocessing}
\label{sec:preprocessing}

\subsection{Loading and Filtering}

To avoid loading the full $\sim$20,000-gene matrix into memory for samples
outside our target cohort, we first read only the header row of the TSV to
identify columns whose 12-character patient barcode prefix appears in the
clinical metadata for our six cancer types.  Only those columns are loaded.
This reduces I/O time significantly for the 1.8\,GB expression file.

\subsection{Log-Transformation and NaN Handling}

Raw FPKM-UQ values are heavy-tailed and span several orders of magnitude.
We apply a $\log_2(x + 1)$ transformation to stabilise variance and make the
distribution approximately Gaussian --- a common preprocessing step in RNA-seq
analysis.  Missing values (NaN) are filled with 0, representing absent/
unexpressed genes.

\subsection{Variance Filtering}

Genes with near-zero variance across all samples ($\sigma^2 < 10^{-6}$) are
removed, as they carry no discriminative signal.  After filtering, the dataset
retains \textbf{20,297 gene features} across 3,607 samples.

\subsection{Train / Test Split}

We hold out 15\% of samples as a final test set using stratified random
splitting (random seed 42), yielding:
\begin{itemize}[leftmargin=2em,itemsep=0pt]
  \item Development set: 3,065 samples (85\%)
  \item Test set: 542 samples (15\%)
\end{itemize}
The test set is never used during model selection or hyperparameter tuning.

\subsection{Standardisation and PCA}

For clustering (Method 1), we fit a \texttt{StandardScaler} on the development
set and apply PCA (50 components) on the scaled development data.  For
supervised methods (Methods 2 \& 3), all preprocessing is performed \emph{inside
each cross-validation fold} on the fold's training partition only, so that no
information from validation or test data contaminates the fitted transformers
(see Section~\ref{sec:methodology}).

% ─────────────────────────────────────────────────────────────────────────────
\section{Methodology and Implementation}
\label{sec:methodology}

All three methods are implemented from scratch in NumPy (Methods 1 \& 2) or
PyTorch (Method 3).  No scikit-learn estimators are used for the ML models;
scikit-learn is used only for evaluation metrics, data splitting, and
preprocessing utilities.

\subsection{Method 1: PCA + K-Means Clustering}
\label{sec:method1}

\paragraph{PCA.}
We implement PCA via truncated SVD on the mean-centred data matrix.
Given $\mathbf{X} \in \mathbb{R}^{n \times p}$ (centred), we compute
$\mathbf{X} = \mathbf{U}\mathbf{S}\mathbf{V}^\top$ and retain the top $k=50$
right singular vectors as principal component directions.  The projection is
$\mathbf{Z} = \mathbf{X}\mathbf{V}_k$.

The explained variance of component $i$ is $\lambda_i = s_i^2 / (n-1)$, and
the explained variance ratio is $r_i = \lambda_i / \sum_j \lambda_j$ (summed
over \emph{all} components, not just the retained $k$, to give a properly
normalised fraction of total variance).  The top 50 PCs collectively explain
\textbf{63.7\%} of total variance (Figure~\ref{fig:pca_variance}).

\paragraph{K-Means.}
We implement Lloyd's K-Means algorithm with $k = 6$ clusters
(matching the number of cancer types), $n\_\text{init} = 10$ random restarts,
and $\text{tol} = 10^{-4}$.  For each restart we:
\begin{enumerate}[leftmargin=2em,itemsep=0pt]
  \item initialise centroids by sampling $k$ data points at random;
  \item alternate between assigning each point to its nearest centroid
        (Euclidean distance) and recomputing centroids as cluster means;
  \item stop when centroid shift $< \text{tol}$ or after 300 iterations.
\end{enumerate}
After convergence, labels are \emph{re-assigned} using the final centroids
before computing inertia (correcting a subtle ordering bug where labels could
refer to pre-update centroids).  The restart with lowest inertia is kept.

\subsection{Method 2: Logistic Regression (One-vs-Rest)}
\label{sec:method2}

We train $C = 6$ binary logistic classifiers in a one-vs-rest (OvR) scheme.
Each binary classifier minimises the $\ell_2$-regularised cross-entropy loss
via mini-batch gradient descent.  For class $c$ with binary labels
$\tilde{y} \in \{0,1\}^n$, weights $\mathbf{w} \in \mathbb{R}^d$, and
bias $b$:

\begin{equation}
  \mathcal{L}(\mathbf{w}, b) =
    -\frac{1}{n}\sum_{i=1}^n \Bigl[\tilde{y}_i \log \hat{y}_i +
      (1-\tilde{y}_i)\log(1-\hat{y}_i)\Bigr]
    + \frac{\lambda}{2n} \|\mathbf{w}\|_2^2
\end{equation}

where $\hat{y}_i = \sigma(\mathbf{x}_i^\top \mathbf{w} + b)$ and $\sigma$ is
the sigmoid function (clipped to $[-250, 250]$ for numerical stability).
Gradient updates are applied with learning rate $\eta = 0.1$ for 1,000
iterations.  Multi-class prediction selects $\arg\max_c \hat{y}_c$ after
normalising OvR probabilities.

\paragraph{Hyperparameter tuning.}
We search $\lambda \in \{0.001, 0.01, 0.1\}$ using 5-fold stratified
cross-validation on the development set.  Inside each fold, \texttt{StandardScaler}
and PCA are fit on the fold's training partition and applied to the validation
partition --- this is critical to prevent data leakage.  The best $\lambda$ is
selected by mean macro-F1 over the five folds.

\subsection{Method 3: Multi-Layer Perceptron (MLP)}
\label{sec:method3}

We implement a feedforward neural network in PyTorch with the architecture:

\begin{equation*}
  \underbrace{50}_{\text{PCA input}} \to
  \underbrace{512}_{\text{Linear+ReLU+BN+Dropout}} \to
  \underbrace{256}_{\text{Linear+ReLU+BN+Dropout}} \to
  \underbrace{6}_{\text{softmax output}}
\end{equation*}

Each hidden layer uses Batch Normalisation and Dropout ($p=0.2$) for
regularisation.  We train with Adam ($\eta = 0.001$, $\beta_1=0.9$,
$\beta_2=0.999$), batch size 64, for 50 epochs, using
\texttt{CrossEntropyLoss}.  Training and evaluation are performed on CPU.

The MLP is evaluated with the same 5-fold stratified CV framework as the LR
model (preprocessing fit inside each fold), then retrained on the full
development set for final test evaluation.

\subsection{Evaluation Metrics}

\begin{itemize}[leftmargin=2em,itemsep=1pt]
  \item \textbf{Clustering:} Adjusted Rand Index (ARI), Normalised Mutual
        Information (NMI), Silhouette score.
  \item \textbf{Classification:} Macro-averaged F1, Precision, Recall;
        per-class confusion matrix; 5-fold CV mean $\pm$ std F1.
\end{itemize}

Macro-averaging weights each class equally regardless of sample count, making
it appropriate for the mild class imbalance in our dataset.

% ─────────────────────────────────────────────────────────────────────────────
\section{Results}
\label{sec:results}

\subsection{PCA Analysis}

Figure~\ref{fig:pca_variance} shows the per-component and cumulative explained
variance.  The first PC alone explains approximately 18\% of variance, with a
characteristic elbow around PC 8--10.  The top 50 PCs together capture
\textbf{63.7\%} of total variance --- sufficient dimensionality reduction for
meaningful downstream analysis while reducing the feature space from 20,297 to
50 dimensions (a $\sim$400$\times$ reduction).

\begin{figure}[H]
  \centering
  \includegraphics[width=\textwidth]{../results/report_figures/fig2_pca_variance.pdf}
  \caption{PCA explained variance on the development set.
           \textit{Left:} per-component variance (scree plot).
           \textit{Right:} cumulative variance; the dashed red line marks the
           63.7\% explained by the top 50 components.}
  \label{fig:pca_variance}
\end{figure}

Figure~\ref{fig:pca_labels} shows the PCA projection coloured by true cancer
type.  Even in the first two PCs --- which together explain only $\sim$23\% of
variance --- the six cancer types form largely non-overlapping clusters.  PRAD
and BRCA show some proximity in PC1, while GBM (brain) is clearly separated
along PC2.

\begin{figure}[H]
  \centering
  \includegraphics[width=0.72\textwidth]{../results/report_figures/fig3_pca_true_labels.pdf}
  \caption{PCA projection (PC1 vs PC2) coloured by true cancer type label.}
  \label{fig:pca_labels}
\end{figure}

\subsection{Method 1: Unsupervised Clustering}

\begin{table}[H]
  \centering
  \caption{K-Means clustering evaluation on the development set.}
  \label{tab:clustering}
  \begin{tabular}{lr}
    \toprule
    Metric & Score \\
    \midrule
    Adjusted Rand Index (ARI) & 0.794 \\
    Normalised Mutual Information (NMI) & 0.849 \\
    Silhouette Score & 0.228 \\
    \bottomrule
  \end{tabular}
\end{table}

Table~\ref{tab:clustering} summarises the unsupervised results.  An ARI of
0.794 (where 0 is random and 1 is perfect) indicates strong but not complete
recovery of the true partitioning.  Figure~\ref{fig:kmeans} compares the
K-Means cluster assignments with the ground truth labels.

\begin{figure}[H]
  \centering
  \includegraphics[width=\textwidth]{../results/report_figures/fig4_kmeans_vs_truth.pdf}
  \caption{PCA projection (PC1 vs PC2).
           \textit{Left:} samples coloured by true cancer type.
           \textit{Right:} samples coloured by K-Means cluster assignment;
           crosses (\textsf{X}) mark cluster centroids.
           The two panels are visually similar, confirming high ARI.}
  \label{fig:kmeans}
\end{figure}

Figure~\ref{fig:purity} shows the cluster purity heatmap.  Five of the six
clusters align almost exclusively with a single cancer type.  The main source
of confusion is GBM, which at only 168 samples is partially merged with another
cluster.  BRCA, KIRC, LUAD, PRAD, and COAD are each captured in a dedicated,
pure cluster.

\begin{figure}[H]
  \centering
  \includegraphics[width=0.72\textwidth]{../results/report_figures/fig5_cluster_purity.pdf}
  \caption{Cluster composition heatmap.  Colour encodes the proportion of each
           cluster belonging to each cancer type; cell numbers show raw counts.
           Most clusters are dominated by a single cancer type.}
  \label{fig:purity}
\end{figure}

The Silhouette score of 0.228 is consistent with published results on
high-dimensional transcriptomic data: studies on TCGA bulk RNA-seq routinely
report Silhouette scores in the 0.15--0.35 range even for well-validated
subtypes, because within-cluster variance in 50-dimensional PCA space is
inherently large relative to between-cluster separation~\cite{weinstein2013}.
ARI and NMI are more informative here because they measure agreement with known
labels directly rather than relying on distance geometry.

\subsection{Methods 2 \& 3: Supervised Classification}

\paragraph{Cross-validation.}
Figure~\ref{fig:cv_folds} shows per-fold macro-F1 for each model.
Both models are remarkably stable: LR achieves $0.9968 \pm 0.0027$ and MLP
achieves $0.9962 \pm 0.0027$ across five folds.  The two models are
statistically indistinguishable in cross-validation.

\begin{figure}[H]
  \centering
  \includegraphics[width=0.75\textwidth]{../results/report_figures/fig6_cv_folds.pdf}
  \caption{5-fold CV macro-F1 for Logistic Regression (blue) and MLP (red).
           Dashed lines show the cross-fold mean for each model.}
  \label{fig:cv_folds}
\end{figure}

\paragraph{Test set evaluation.}
Table~\ref{tab:classification} reports test-set performance for both classifiers.
Figure~\ref{fig:lr_cm} and Figure~\ref{fig:mlp_eval} show the confusion
matrices; per-class performance is analysed in detail below.

\begin{table}[H]
  \centering
  \caption{Test set classification performance (542 samples, 15\% hold-out).
           LR hyperparameter: $\lambda=0.001$.}
  \label{tab:classification}
  \begin{tabular}{lcccc}
    \toprule
    Model & CV F1 (mean $\pm$ std) & Test F1 & Test Precision & Test Recall \\
    \midrule
    Logistic Regression & $0.9968 \pm 0.0027$ & 0.9927 & 0.9877 & 0.9982 \\
    MLP [512, 256]      & $0.9962 \pm 0.0027$ & 0.9922 & 0.9877 & 0.9973 \\
    \bottomrule
  \end{tabular}
\end{table}

\begin{figure}[H]
  \centering
  \includegraphics[width=0.62\textwidth]{../results/report_figures/fig7_lr_confusion.pdf}
  \caption{Logistic Regression confusion matrix on the held-out test set.
           Off-diagonal entries are concentrated in the GBM row.}
  \label{fig:lr_cm}
\end{figure}

\begin{figure}[H]
  \centering
  \includegraphics[width=\textwidth]{../results/report_figures/fig8_mlp_eval.pdf}
  \caption{\textit{Left:} MLP confusion matrix on the test set.
           \textit{Right:} MLP training dynamics --- loss decreases smoothly and
           validation accuracy converges to $\approx 1.0$ within $\sim$20
           epochs.}
  \label{fig:mlp_eval}
\end{figure}

\paragraph{Per-class analysis.}
Figure~\ref{fig:perclass} shows per-class binary F1 for both models.  All
cancer types except GBM achieve F1 $= 1.00$ for both classifiers.  GBM
obtains F1 $= 0.96$, consistent with its smaller sample size (168 total,
$\approx 25$ in the test set) and partial cluster overlap observed in the
K-Means analysis.

\begin{figure}[H]
  \centering
  \includegraphics[width=0.78\textwidth]{../results/report_figures/fig9_perclass_f1.pdf}
  \caption{Per-class F1 score on the test set.  GBM is the only class where
           either model makes errors; all other types are classified perfectly.}
  \label{fig:perclass}
\end{figure}

\paragraph{LR vs MLP comparison.}
Figure~\ref{fig:summary} provides a direct head-to-head comparison of macro
F1, Precision, and Recall on the test set.  The MLP offers a negligible
$-0.05$ percentage point difference in macro-F1 compared with logistic
regression, confirming that the classification problem is essentially linearly
separable in the 50-dimensional PCA space.

\begin{figure}[H]
  \centering
  \includegraphics[width=0.72\textwidth]{../results/report_figures/fig10_summary_comparison.pdf}
  \caption{Test set macro-averaged metrics for both supervised classifiers.
           LR and MLP perform equivalently across all three metrics.}
  \label{fig:summary}
\end{figure}

% ─────────────────────────────────────────────────────────────────────────────
\section{Analysis and Discussion}
\label{sec:discussion}

\subsection{Why Are Accuracies So High?}

The near-perfect classification results ($>$99\% macro-F1) initially appear
surprising but are consistent with the biology and prior literature:

\begin{itemize}[leftmargin=2em,itemsep=2pt]
  \item \textbf{Tissue-of-origin dominates expression.}  The six cancer types
        originate from entirely different tissues (breast, lung, kidney,
        prostate, colon, brain).  Each tissue has a distinct baseline
        transcriptome, and cancer largely preserves these tissue-specific
        signatures~\cite{weinstein2013}.  The signal separating cancer types
        is therefore vastly larger than the within-type tumour heterogeneity.

  \item \textbf{Batch correction removes technical noise.}  The PANCAN dataset
        uses EBPlusPlus batch correction, which removes sequencing artefacts
        that would otherwise inflate or deflate expression estimates across
        batches.  The remaining variance is predominantly biological.

  \item \textbf{Rich feature space.}  Even after $\log_2$ transformation and
        PCA reduction, 50 components capturing 63.7\% of total variance
        provide abundant discriminative signal for a 6-class problem.

  \item \textbf{Literature precedent.}  Published benchmarks on TCGA Pan-Cancer
        classification report 90--99\% accuracy using RNA-seq
        features~\cite{weinstein2013,hutter2018}.  Our results are within the
        expected range.
\end{itemize}

The task solved here is effectively ``which organ does this tumour come from?''
--- a much easier problem than, e.g., distinguishing PAM50 molecular subtypes
within breast cancer, or predicting treatment response from expression data.

\subsection{Why Does MLP Offer No Gain Over LR?}

The near-zero performance gap ($-0.05$ pp) between LR and MLP is a direct
consequence of \emph{linear separability} in PCA space.  The PCA scatter plot
(Figure~\ref{fig:pca_labels}) already shows clean cluster separation in just
two dimensions.  If the decision boundary is linear, a non-linear model adds
capacity without improving accuracy --- and may even slightly underperform due
to the stochastic nature of mini-batch training.  This confirms the proposal's
prediction: ``we will compare Methods 2 and 3 directly to assess the gain from
non-linearity.''

\subsection{Unsupervised vs Supervised Performance}

The ARI of 0.794 for K-Means, while high, trails the supervised classifiers
by a substantial margin.  This gap is expected: K-Means has no access to labels
and must discover structure purely from geometry.  As shown in
Figure~\ref{fig:purity}, five of the six clusters map cleanly to a single
cancer type; the residual error comes almost entirely from the GBM cluster,
as discussed in Section~\ref{sec:results}.

\subsection{Data Leakage Prevention}

A critical methodological concern in machine learning pipelines on genomic data
is \emph{data leakage} --- allowing information from the validation or test set
to influence preprocessing fitted on training data.  In this work, all
\texttt{StandardScaler} and PCA fits are performed \emph{inside each CV fold}
on the fold's training partition only.  The held-out test set is never seen
until the final evaluation.  This design ensures that the reported metrics are
faithful estimates of generalisation performance.

\subsection{Limitations}

\begin{itemize}[leftmargin=2em,itemsep=1pt]
  \item \textbf{Sample imbalance.}  GBM's lower F1 (0.96 vs 1.00 for all
        other classes) traces directly to its smaller cohort (168 samples,
        $7\times$ fewer than BRCA); future work could address this with
        data augmentation or over-sampling strategies.
  \item \textbf{Task difficulty.}  Classifying across tissue-of-origin is a
        relatively easy task.  Harder biologically relevant tasks (subtype
        discovery within a cancer type, survival prediction) would provide a
        more discriminating benchmark.
  \item \textbf{Architecture search.}  The MLP architecture (512, 256) was
        fixed; a systematic hyperparameter search could potentially improve
        performance on harder tasks.
  \item \textbf{Interpretability.}  The logistic regression weights
        $\mathbf{w} \in \mathbb{R}^{50}$ operate in PCA space rather than gene
        space, making direct biological interpretation impossible.  The natural
        path forward is to project these weights back to gene space via
        $\tilde{\mathbf{w}} = \mathbf{V}_k \mathbf{w}$, where $\mathbf{V}_k$
        are the top-$k$ PCA loadings; the resulting $\tilde{\mathbf{w}} \in
        \mathbb{R}^{20297}$ ranks each gene's contribution to each OvR
        classifier, enabling pathway enrichment analysis and comparison with
        known tissue-specific markers.
\end{itemize}

% ─────────────────────────────────────────────────────────────────────────────
\section{Conclusions}
\label{sec:conclusions}

We implemented and evaluated a complete machine learning pipeline for cancer
type discovery and classification from TCGA Pan-Cancer RNA-seq data:

\begin{enumerate}[leftmargin=2em,itemsep=2pt]
  \item \textbf{PCA + K-Means} (unsupervised) recovers the six cancer types
        with ARI = 0.794 and NMI = 0.849, demonstrating that bulk RNA-seq
        contains strong intrinsic structure aligned with tissue-of-origin.

  \item \textbf{Logistic Regression with $\ell_2$ regularisation} achieves
        macro-F1 = 0.9927 on the held-out test set (5-fold CV: $0.9968 \pm
        0.0027$), confirming that the six-class problem is nearly linearly
        separable in 50-dimensional PCA space.

  \item \textbf{MLP} (PyTorch, two hidden layers) achieves macro-F1 = 0.9922
        --- statistically equivalent to logistic regression ($-0.05$ pp gap)
        --- indicating that the added non-linearity provides no benefit when
        the decision boundaries are already approximately linear.
\end{enumerate}

The consistently high performance across all three methods underscores that the
fundamental biological signal (tissue-of-origin) is robustly encoded in bulk
gene expression, and that even simple, properly implemented algorithms can
extract it with high fidelity when applied to well-curated, batch-corrected
data.

% ─────────────────────────────────────────────────────────────────────────────
\section*{Implementation Notes}

All code is implemented in Python 3.  PCA and K-Means are written in NumPy;
Logistic Regression is implemented from scratch in NumPy; the MLP uses PyTorch
1.x.  Data download and preprocessing use \texttt{requests}, \texttt{pandas},
and \texttt{numpy}.  Evaluation metrics are computed with \texttt{scikit-learn}.
The full source code and results are available in the project repository.

\begin{itemize}[leftmargin=2em,itemsep=1pt]
  \item \texttt{src/data\_loader.py} --- GDC API download, log transform, variance filter
  \item \texttt{src/pca.py} --- PCA via truncated SVD (NumPy)
  \item \texttt{src/kmeans.py} --- Lloyd's K-Means with multiple restarts (NumPy)
  \item \texttt{src/logistic\_regression.py} --- OvR logistic regression with $\ell_2$ (NumPy)
  \item \texttt{src/mlp.py} --- feedforward MLP (PyTorch)
  \item \texttt{src/evaluate.py} --- metrics and plotting utilities
  \item \texttt{src/main.py} --- end-to-end pipeline with 5-fold CV
  \item \texttt{src/generate\_report\_figures.py} --- report-quality figure generation
\end{itemize}

% ─────────────────────────────────────────────────────────────────────────────
\bibliographystyle{plainnat}
\begin{thebibliography}{9}

\bibitem{weinstein2013}
Weinstein, J.~N., Collisson, E.~A., Mills, G.~B., et al.\ (2013).
The Cancer Genome Atlas Pan-Cancer analysis project.
\textit{Nature Genetics}, 45(10), 1113--1120.

\bibitem{hutter2018}
Hutter, C., \& Zenklusen, J.~C.\ (2018).
The Cancer Genome Atlas: creating lasting value beyond its data.
\textit{Cell}, 173(2), 283--285.

\end{thebibliography}

\end{document}
